% schema: latex-v1
\documentclass[11pt]{article}
\usepackage{cslv3}
\usepackage[margin=1in]{geometry}
\usepackage{xunicode}
\usepackage{xltxtra}
\title%!(MISSING CLOSE BRACE)s}
%!(EXTRA C5\_bridge\_mode)\author{CSLv3 auto-emit}
\date{\today}

\begin{document}
\maketitle
\tableofcontents

\section%!(MISSING CLOSE BRACE)s}\label%!(MISSING CLOSE BRACE)s}
%!(EXTRA bridge.example, sec:bridgeexample)\section%!(MISSING CLOSE BRACE)s}\label%!(MISSING CLOSE BRACE)s}
%!(EXTRA GOAL, sec:goal)\begin{lstlisting}[language=csl]
# tests: mixed EN prose + CSL fragments (§09 bridge mode)
# glyph-class: Comment preservation, ident-heavy lines, embedded structured blocks
# features: prose coexisting with parseable CSLv3 in the same document
# This document demonstrates bridge mode where English prose sits next to
# structured CSLv3. Readers unfamiliar with the notation can still follow
# the intent while the parser extracts only the structured regions.
# A short English paragraph would live here in practice. The compiler
# treats comment lines as prose and ignores them.
summary : str
\end{lstlisting}
\section%!(MISSING CLOSE BRACE)s}\label%!(MISSING CLOSE BRACE)s}
%!(EXTRA STRUCTURED, sec:structured)\begin{lstlisting}[language=csl]
def Request't ⟨
  method : str
  path : str
  body : str
⟩
\end{lstlisting}
\begin{lstlisting}[language=csl]
def Response't ⟨
  status : u16
  body : str
⟩
\end{lstlisting}
\section%!(MISSING CLOSE BRACE)s}\label%!(MISSING CLOSE BRACE)s}
%!(EXTRA OPS, sec:ops)\begin{lstlisting}[language=csl]
# Inline explanation: the next function's contract is that any valid
# Request must map to exactly one Response. English above; CSL below.
fn handle (req : &Request't) -> Response't = Response't
\end{lstlisting}
\section%!(MISSING CLOSE BRACE)s}\label%!(MISSING CLOSE BRACE)s}
%!(EXTRA INVARIANTS, sec:invariants)\begin{lstlisting}[language=csl]
⌈determinism⌉
\end{lstlisting}
\begin{lstlisting}[language=csl]
⌈request-maps-to-one-response⌉
\end{lstlisting}

\end{document}
